\begin{textsample}{POS Dim 4 – 2014\_c4 – Score 1.00 – 2014\_c4/enade\_2014\_letras\_por\_ing\_116500ce32.txt}  \label{ex:f4_pos_208}
CAPÍTULO 4 PERCEPÇÃO DA PROVA As análises feitas neste capítulo tratam das percepções dos concluintes da Área de Letras - Português e Inglês ( Licenciatura ) sobre a prova aplicada no ENADE/2014.
Estas percepções foram mensuradas por meio de nove questões que avaliaram desde o grau de dificuldade da prova até o tempo gasto para concluí -la.
As percepções sobre a prova foram relacionadas com o desempenho dos estudantes e com a Grande Região de funcionamento do curso.
O questionário de percepção da prova encontra -se ao final do Anexo IV, que traz a reprodução da prova.
O desempenho dos estudantes foi classificado em quatro quartos.
Para tanto, esse desempenho foi ordenado de forma ascendente.
O percentil 25, P25, também conhecido como primeiro quartil, é a nota de desempenho que deixa um quarto ( 25 \% ) dos valores observados abaixo e três quartos acima.
A Figura 1 apresenta uma ilustração deste conceito.
O quarto inferior de desempenho é composto pelas notas abaixo do primeiro quartil.
Já o percentil 75, P75, também conhecido como terceiro quartil, é o valor para o qual há três quartos ( 75 \% ) dos dados abaixo e um quarto acima dele.
O quarto superior de desempenho é composto pelas notas iguais ou acima do terceiro quartil.
O percentil 50, P50, também conhecido como mediana, é o valor que divide as notas em dois conjuntos de igual tamanho.
O segundo quarto inclui valores entre o primeiro quartil ( P25 ) e a mediana.
O terceiro quarto contém os valores entre a mediana ( P50 ) e o terceiro quartil ( P75 ).
Vale ressaltar que percentis, quartis e medianas são pontos que não obrigatoriamente pertencem ao conjunto original de dados, ao passo que os quartos são subconjuntos dos dados originais.
A seguir, serão apresentados gráficos com resultados selecionados, relativos às nove questões avaliadas por grupos de estudantes.
Os gráficos apresentam nas barras o percentual de alunos que assinalaram uma das opções ou a soma das porcentagens daqueles que assinalaram duas ( ou três ) delas.
Por exemplo, para as questões 1 e 2, os gráficos apresentam a porcentagem total de participantes que assinalaram as opções ( D ) difícil e ( E ) muito difícil.
Em cada barra foram assinalados também os extremos do intervalo de confiança de 95 \% como linhas verticais unidas por uma linha horizontal na forma da letra H maiúscula.
As Tabelas no Anexo II apresentam os valores absolutos e a distribuição percentual21 das alternativas válidas das nove questões, segundo o mesmo recorte de desempenho dos alunos e Grande Região de funcionamento do curso. 4.1 GRAU DE DIFICULDADE DA PROVA 4.1.1 Componente de Formação Geral Ao avaliarem “ Qual o grau de dificuldade desta prova na parte de Formação Geral? ” ( Questão 1 ), 33,0 \% do grupo de inscritos e presentes optaram pelas alternativas difícil ou muito difícil.
Entretanto, para mais da metade dos estudantes ( 57,6 \% ), o Componente de Formação Geral da prova foi considerado com grau de dificuldade médio ( Gráfico 4.1, Gráfico 4.2 e, no Anexo II, a Tabela II.1 ).
O percentual de estudantes que consideraram a prova como difícil ou muito difícil foi maior na região Norte, onde a proporção foi de 37,9 \%, enquanto a de menor incidência foi a Sudeste, com 30,4 \%.
No Gráfico 4.1, é possível observar que esta diferença é estatisticamente significativa.
Nas Grandes Regiões, a proporção de presentes à prova que consideraram o Componente de Formação Geral como sendo de grau de dificuldade médio esteve entre 54,5 \% na região Norte e 59,4 \% na região Sudeste. 21 Cumpre lembrar uma das convenções para tabelas numéricas ( pág. iii ) sobre a possibilidade da soma das partes não resultar em 100 \% por questões de arredondamento.
O percentual de alunos que consideraram a prova difícil ou muito difícil foi decrescente em função dos quartos de desempenho : 42,5 \% no primeiro quarto e 22,4 \% no quarto, grupo de melhor desempenho na prova.
Nos quartos de desempenho intermediários, a proporção de alunos que consideraram a prova difícil ou muito difícil foi de 37,1 \% no 2º quarto e 30,5 \% no 3º quarto.
As diferenças entre os quartos são estatisticamente significativas.
Para todos os quartos de desempenho, a alternativa modal para esta pergunta foi médio, com 50,4 \% e 62,5 \% dos respondentes nos quartos extremos, primeiro e quarto, respectivamente, valores crescentes com o desempenho. 4.1.2 Componente de Conhecimento Específico Ao responderem à Questão 2 – “ Qual o grau de dificuldade desta prova na parte de Componente Específico? ” – 36,7 \% do grupo de estudantes classificaram -na como difícil ou muito difícil.
Além disso, o Componente de Conhecimento Específico da prova foi considerado com grau de dificuldade médio por 56,0 \% dos alunos ( Gráfico 4.3, Gráfico 4.4, e, no Anexo II, a Tabela II.2 ).
A análise das respostas dos estudantes quanto ao grau de dificuldade do Componente de Conhecimento Específico da prova, agregado por Grande Região, mostra que a diferença entre a maior e as três menores proporções de alunos que a avaliaram como difícil ou muito difícil é estatisticamente significativa : a maior na região Centro-Oeste ( 41,0 \% ) e as menores proporções nas regiões Sudeste ( 35,3 \% ), Sul ( 35,6 \% ) e Nordeste ( 36,3 \% ).
O percentual de alunos que classificaram o grau de dificuldade como médio, no Componente de Conhecimento Específico, variou de 53,5 \% a 57,3 \%, para as regiões Norte e Sudeste, respectivamente.
A proporção dos que classificaram a parte específica como difícil ou muito difícil variou de 29,5 \% ( 4º quarto ) a 41,6 \% ( 1º quarto ), valores decrescentes com o desempenho.
Considerando -se essas opiniões sobre a avaliação da dificuldade das questões do Componente de Conhecimento Específico da prova, de acordo com o desempenho dos estudantes, observa -se que há diferença estatisticamente significativa do resultado do quarto de melhor desempenho em relação aos demais quartos.
A alternativa modal para a Questão 2 foi o grau médio, com 51,0 \% do quarto inferior e 61,1 \% do superior optando por esta resposta, valores crescentes com o desempenho. 4.2 EXTENSÃO DA PROVA EM RELAÇÃO AO TEMPO TOTAL Indagados quanto à extensão da prova, em relação ao tempo total oferecido para a sua resolução ( Questão 3 ), os estudantes apontaram, com maior incidência, a alternativa que considerava a extensão adequada, para todas as agregações consideradas ( Gráfico 4.5, Gráfico 4.6, e, no Anexo II, a Tabela II.3 ).
O percentual de alunos que responderam ser a extensão da prova adequada foi de 42,1 \%.
Já 51,1 \% dos inscritos presentes consideraram que o exame foi longo ou muito longo e 6,8 \% o avaliaram como curto ou muito curto.
Entre as Grandes Regiões a proporção daqueles que avaliaram a prova como longa ou muito longa em relação ao tempo total destinado à sua resolução variou pouco : de 45,7 \% na região Norte até 53,6 \% na região Sudeste.
A diferença entre as regiões Norte e Sudeste é estatisticamente significativa.
Considerando -se o desempenho dos alunos, nota -se ainda que no quarto de desempenho inferior, 40,9 \% consideraram a extensão da prova adequada e no de melhor desempenho ( quarto superior ), 42,2 \%.
Nos quartos intermediários, esta proporção foi 41,5 \% no segundo quarto e 43,7 \% no terceiro.
No Gráfico 4.6, pode -se constatar que, grosso modo, há uma tendência crescente da proporção de estudantes que consideraram a prova longa ou muito longa em função dos quartos de desempenho.
Observa -se que não há diferença estatisticamente significativa entre as proporções de escolha desta resposta no primeiro e no último quarto. 4.3 COMPREENSÃO DOS ENUNCIADOS DAS QUESTÕES 4.3.1 Componente de Formação Geral Com relação aos enunciados das questões do Componente de Formação Geral ( Questão 4 ), as opiniões foram positivas, já que 75,7 \% dos alunos avaliados consideraram os enunciados de todas ou da maioria das questões claros e objetivos ( Gráfico 4.7, Gráfico 4.8, e, no Anexo II, a Tabela II.4 ).
Na análise regional, a percentagem de estudantes que avaliaram que todos ou a maioria dos enunciados das questões do Componente de Formação Geral estavam claros e objetivos variou de 70,3 \% na região Nordeste a 81,0 \% na região Sul, sendo esta diferença estatisticamente significativa.
A análise das percepções dos estudantes sobre a clareza e objetividade dos enunciados permite afirmar que todos, ou a maioria dos enunciados de questões relativas ao Componente de Formação Geral, foram considerados claros e objetivos para a maior parte dos respondentes ( maior do que 70 \% em todas as regiões e maior do que 62 \% para todos os quartos de desempenho ).
Segundo o desempenho, observa -se que a proporção dos que emitiram esta opinião cresce conforme o desempenho aumenta, com diferenças estatisticamente significativas entre toso os quartos de desempenho.
No quarto superior, a clareza e objetividade de todos ou da maioria dos enunciados das questões foi percebida por 86,1 \% dos alunos.
No quarto inferior esta percentagem foi 62,7 \%. 4.3.2 Componente de Conhecimento Específico Com relação aos enunciados das questões do Componente de Conhecimento Específico da prova, para 74,6 \% dos estudantes avaliados da Área de Letras - Português e Inglês ( Licenciatura ), a clareza e a objetividade ( Questão 5 ) estavam presentes em todas ou na maioria das questões ( Gráfico 4.9, Gráfico 4.10, e no Anexo II, a Tabela II.5 ).
A maioria dos estudantes de todas as Grandes Regiões brasileiras considerou claros e objetivos todos ou a maioria dos enunciados das questões do Componente de Conhecimento Específico da prova, percentual sempre maior ou igual a 68,6 \%.
A diferença entre a maior e a menor proporção, 78,8 \% ( Sul ) e 68,6 \% ( Nordeste ), é estatisticamente significativa.
A proporção de estudantes que consideraram os enunciados das questões claros e objetivos apresenta uma tendência crescente em relação ao aumento de desempenho : mais elevada no quarto superior ( 83,6 \% ) se comparada ao quarto inferior de desempenho ( 63,4 \% ).
As diferenças entre todos os quartos são estatisticamente significativas. 4.4 SUFICIÊNCIA DAS INFORMAÇÕES/INSTRUÇÕES FORNECIDAS Ao avaliarem as informações/instruções fornecidas para a resolução das questões ( Questão 6 ), 83,1 \% dos respondentes da Área de Letras - Português e Inglês ( Licenciatura ) de todo o Brasil afirmaram que estas eram até excessivas ou suficientes em todas ou na maioria das questões ( Gráfico 4.11, Gráfico 4.12, e, no Anexo II, a Tabela II.6 ).
Quanto à distribuição de respondentes pelas Grandes Regiões observa -se que a proporção de estudantes que consideraram as informações/instruções fornecidas até excessivas ou suficientes em todas ou na maioria das questões foi sempre superior a 77,0 \%, chegando a 89,4 \% na região Sul.
A diferença entre o maior percentual, obtido na região Sul ( 89,4 \% ), e o menor, da região Nordeste ( 77,2 \% ) é estatisticamente significativa.
O percentual de participantes que avaliaram as informações/instruções como até excessivas ou suficientes em todas ou na maioria das questões foi mais elevado no quarto superior ( 90,7 \% ), percentual superior à média nacional ( 83,1 \% ).
Já no quarto inferior, a suficiência das informações/instruções declarada como até excessiva em todas ou na maioria das questões foi percebida por 72,3 \% dos respondentes.
Estas proporções são crescentes com os quartos de desempenho.
Notam -se diferenças estatisticamente significativas entre as opiniões de estudantes de todos os quartos de desempenho, como mostra o Gráfico 4.12. 4.5 DIFICULDADE ENCONTRADA AO RESPONDER À PROVA Perguntados sobre as dificuldades com as quais se depararam ao responder à prova ( Questão 7 ), 16,6 \% dos estudantes apontaram o desconhecimento do conteúdo.
Para 41,5 \%, a forma diferente de abordagem do conteúdo foi indicada como dificuldade.
Já a falta de motivação para fazer a prova foi a dificuldade apontada por 16,5 \% dos respondentes.
Considerando -se todo o Brasil, 18,3 \% dos respondentes afirmaram que não tiveram qualquer tipo de dificuldade para responder à prova ( Tabela II.7 no Anexo II ).
Os Gráficos 4.13 e 4.14 apresentam os percentuais de estudantes que apontaram o desconhecimento do conteúdo como dificuldade percebida ao responder à prova.
Na análise por Grandes Regiões, o percentual de inscritos e presentes que apontaram o desconhecimento do conteúdo como dificuldade ao responder à prova não superou 23,0 \%.
Os percentuais variaram de 14,4 \% na região Sudeste a 23,0 \% na Norte, sendo esta diferença significativa estatisticamente.
A forma diferente de abordagem do conteúdo foi a escolha modal dos estudantes, com percentuais que variaram de 35,7 \% ( região Sul ) a 44,2 \% ( Centro-Oeste ).
O percentual de alunos que citaram a falta de motivação como dificuldade variou de 12,7 \% ( região Norte ) a 19,6 \% ( região Nordeste ).
Os que declararam não ter qualquer dificuldade para responder à prova variaram de 11,8 \% na região Norte a 21,1 \% na Sul.
Com relação aos quartos de desempenho, o desconhecimento do conteúdo foi a opção escolhida por 14,5 \% dos estudantes do quarto superior e 17,9 \% do quarto inferior.
A alternativa modal para os alunos, quando agregados pelos quartos de desempenho, foi que a dificuldade encontrada foi causada pela forma diferente de abordagem do conteúdo : 44,5 \% no quarto inferior e 36,4 \% do quarto superior assim o responderam, valores decrescentes com o desempenho. 4.6 CONTEÚDOS DAS QUESTÕES OBJETIVAS DA PROVA Ao analisarem os conteúdos das questões objetivas da prova ( Questão 8 ), um percentual muito pequeno dos estudantes avaliados, apenas 7,1 \%, afirmou que não estudou ainda a maioria desses conteúdos ( Gráficos 4.15, Gráfico 4.16, e a Tabelas II.8 no Anexo II ).
A maioria ( 66,4 \% ) afirmou ter estudado e aprendido muitos ou todos os conteúdos.
Na análise por Grande Região, a proporção de respondentes que escolheram a opção não estudou ainda a maioria desses conteúdos foi também pequena.
Observa -se que nas regiões Norte ( 12,0 \% ), Centro-Oeste ( 9,3 \% ) e Nordeste ( 8,6 \% ), apesar de pequenas, as proporções foram maiores do que a média nacional ( 7,1 \% ). pequenas, as proporções foram maiores do que a média nacional ( 7,1 \% ).
Observa -se diferença estatisticamente significativa entre algumas das regiões.
Em todas as regiões, a maioria dos presentes afirmou ter estudado e aprendido muitos ou todos os conteúdos, com proporções variando entre 53,3 \% na região Norte e 75,5 \% na Sul.
Considerando -se separadamente as opiniões de estudantes dos quatro quartos de desempenho, observa -se que, no quarto inferior, 12,8 \% ofereceram como resposta que não estudou ainda a maioria desses conteúdos, sendo 2,7 \% os do quarto superior com a mesma resposta.
A diferença entre os alunos que optaram por este motivo de dificuldade nos quartos extremos é estatisticamente significativa, e os valores são decrescentes com o desempenho.
Tendo em conta o quarto superior, 84,9 \% dos alunos afirmaram ter estudado e aprendido muitos ou todos os conteúdos.
No outro extremo, no primeiro quarto, 47,4 \% optaram pelas mesmas categorias, valores crescentes com o desempenho.. 4.7 TEMPO GASTO PARA CONCLUIR A PROVA Ao responderem sobre o tempo de conclusão da prova ( Questão 9 ), mais da metade dos estudantes ( 72,8 \% ) afirmou ter gasto entre duas e quatro horas ( Gráfico 4.17, Gráfico 4.18 e, no Anexo II, a Tabela II.9 ).
Considerando -se as cinco Grandes Regiões brasileiras, o percentual dos que utilizaram entre duas e quatro horas para finalizar a prova nas regiões Sul ( 70,2 \% ), Sudeste ( 72,1 \% ) e Norte ( 72,6 \% ), foi inferior ao percentual nacional.
Nas demais Grandes Regiões o percentual de alunos que dispensaram entre duas e quatro horas para concluir a prova, não \textbf{ultrapassou} 75,0 \%, como mostra o Gráfico 4.17.
Uma vez sendo analisadas as alternativas escolhidas pelos estudantes que se situam nos diferentes quartos de desempenho, observa -se uma tendência crescente : uma maior proporção de participantes no quarto superior declarou ter gasto entre duas e quatro horas para concluir a prova quando comparadas com os dos quartos inferiores.
As diferenças entre o primeiro ( 65,5 \% ) e o segundo ( 72,6 \% ), bem como entre estes e o quarto de maior desempenho ( 79,1 \% ) são estatisticamente significativas.

% matched lemmas: ultrapassar
\end{textsample}
